\section{The next question: balanced summands}\label{sec:balanced}
Existence still leaves the shape of a decomposition largely open. It is
natural to ask for the two summands to have comparable size.
The supplementary NumPy verifier for $\AA+\AA$ also tests $n/3\le a\le b\le2n/3$. In the entire range
$2\le n<100{,}000{,}000$, $n\equiv2\pmod6$, precisely $8$ and $44$
fail. This motivates, but does not prove, the sharper statement:
\begin{conjecture}
Every positive integer $n\equiv2\pmod6$ other than $8$ and $44$ admits $n=a+b$ with
$a,b\in\AA$ and $n/3\le a,b\le2n/3$.
\end{conjecture}
The present counting argument leaves the two norm sizes free. To settle
this last question, we would need control of the ratio
$N(z)/(N(z)+N(w))$, or a different construction that keeps the two values
in the desired interval. The finite search gives a clear target, but it
does not yet provide that extra argument.
